% =============================================================
% Capitolo: La Seller Pipeline — Analisi della Reputazione del Venditore
% Progetto: MCP_ECOM — Agente e-commerce basato su MCP
% =============================================================

\chapter{La Seller Pipeline: Analisi della Reputazione del Venditore}
\label{chap:seller-pipeline}

% ─────────────────────────────────────────────────────────────
\section{Motivazione e Posizione nel Sistema}
\label{sec:seller-intro}
% ─────────────────────────────────────────────────────────────

In un marketplace come eBay, la scelta del prodotto giusto è solo metà
della decisione d'acquisto: l'altra metà riguarda la \emph{fiducia nel
venditore}. Un articolo con prezzo competitivo può nascondere rischi
reali se il venditore ha una storia di feedback negativi, spedizioni
lente o gestione controversie carente.

La \textbf{Seller Pipeline} è il sottosistema del progetto MCP\_ECOM
dedicato alla valutazione automatica e quantitativa della reputazione
di un venditore eBay. Viene invocata dal tool MCP
\texttt{analyze\_seller} e produce, come output, un insieme strutturato
di metriche — \emph{trust score}, \emph{sentiment score}, metadati
anagrafici e un sommario generato da LLM — che l'agente può
presentare direttamente all'utente o usare come segnale di ranking nella
pipeline di ricerca prodotti.

\begin{figure}[htbp]
    \centering
    \fbox{\parbox{0.82\linewidth}{\centering\vspace{2cm}
        \textit{[Diagramma: MCP Tool analyze\_seller → Seller Pipeline
        → (eBay Trading API $\parallel$ Redis Cache) → NLP Sentiment
        → Trust Score → AI Summary → Output normalizzato]}\\
        \vspace{2cm}}}
    \caption{Flusso ad alto livello della Seller Pipeline.}
    \label{fig:seller-pipeline-flow}
\end{figure}


% ─────────────────────────────────────────────────────────────
\section{Architettura della Pipeline}
\label{sec:seller-arch}
% ─────────────────────────────────────────────────────────────

Il punto di entrata è la funzione \texttt{run\_seller\_pipeline()}
(\texttt{app/services/seller\_pipeline.py}). Il flusso si articola in
cinque fasi principali:

\begin{enumerate}
    \item \textbf{Cache check} — ricerca in Redis del risultato già
        calcolato per la coppia \texttt{(seller\_name, limit)}.
    \item \textbf{Fetch parallelo} — acquisizione asincrona e simultanea
        di feedback, dati anagrafici e informazioni sul negozio
        dall'API Trading di eBay.
    \item \textbf{Analisi NLP del sentiment} — calcolo del punteggio
        di sentiment su ciascun feedback testuale tramite embedding
        semantici.
    \item \textbf{Calcolo del trust score} — aggregazione di feedback
        ratio, sentiment score, recency decay e confidence in un unico
        indice composito.
    \item \textbf{AI Summary} — generazione di un sommario testuale
        sintetico tramite LLM, in lingua italiana, basato su un
        campione dei feedback reali.
\end{enumerate}

Il risultato viene infine serializzato, salvato su Redis e restituito
al chiamante.

\subsection{Parallelismo Asincrono}
\label{subsec:seller-async}

Le tre chiamate verso l'API eBay — \texttt{GetFeedback},
\texttt{GetUser} e \texttt{GetStore} — vengono lanciate in parallelo
tramite \texttt{asyncio.gather()}, eliminando la latenza sequenziale
che si accumulerebbe attendendo ciascuna risposta prima di avviare la
successiva:

\begin{verbatim}
feedbacks, user_details, store_details = await asyncio.gather(
    get_seller_feedback(seller_name, limit=needed),
    get_user_details(seller_name),
    get_store_details(seller_name),
)
\end{verbatim}

Analogamente, le due operazioni CPU-intensive — calcolo del sentiment
e calcolo del trust score — vengono eseguite nel thread pool di default
tramite \texttt{asyncio.to\_thread()}, evitando il blocco dell'event
loop durante l'inferenza del modello di embedding.


% ─────────────────────────────────────────────────────────────
\section{Acquisizione Dati: eBay Trading API}
\label{sec:seller-data}
% ─────────────────────────────────────────────────────────────

L'acquisizione dei dati avviene tramite l'eBay \textbf{Trading API}
(SOAP/XML), autenticata con token OAuth seller-level. Il modulo
\texttt{app/services/feedback.py} gestisce tre chiamate distinte.

\subsection{GetFeedback — Fetch Multi-Pagina Asincrono}
\label{subsec:seller-feedback}

La chiamata \texttt{GetFeedback} recupera i feedback lasciati dagli
acquirenti sul venditore. Dato che l'API eBay pagina i risultati a
blocchi di 50 elementi, il sistema calcola automaticamente il numero
di pagine necessarie a soddisfare il \texttt{limit} richiesto
e lancia \emph{tutte le pagine in parallelo}:

\begin{verbatim}
tasks = [
    fetch_feedback_page(client, username, page, per_page)
    for page in range(1, max_pages + 1)
]
results = await asyncio.gather(*tasks)
\end{verbatim}

Per ciascun feedback viene eseguita la conversione della valutazione
qualitativa eBay (\texttt{"Positive"}, \texttt{"Neutral"},
\texttt{"Negative"}) in un valore numerico su scala 1–5,
compatibile con la pipeline di scoring:

\begin{table}[htbp]
\centering
\caption{Mapping valutazione eBay → valore numerico}
\label{tab:rating-map}
\begin{tabular}{cc}
\hline
\textbf{eBay CommentType} & \textbf{Valore numerico} \\
\hline
Positive  & 5 \\
Neutral   & 3 \\
Negative  & 1 \\
\hline
\end{tabular}
\end{table}

Il parser XML include protezione contro l'XML injection via
\texttt{xml.sax.saxutils.escape()} e gestione degli errori eBay
(\texttt{Ack = Failure / PartialFailure}) prima del parsing dei nodi
dati, prevenendo l'ingestione silenziosa di risposte di errore.

\subsection{GetUser e GetStore — Metadati Anagrafici}
\label{subsec:seller-meta}

In parallelo al fetch dei feedback, vengono recuperati:

\begin{itemize}
    \item \textbf{GetUser}: data di registrazione, paese, feedback
        score globale eBay, stato dell'account.
    \item \textbf{GetStore}: nome del negozio eBay, URL del logo,
        descrizione editoriale dello store.
\end{itemize}

Questi metadati non concorrono al calcolo del trust score, ma
arricchiscono il profilo presentato all'utente e sono utili
all'agente per contestualizzare la raccomandazione (es.\ un venditore
registrato da 10 anni con negozio dedicato ispira più fiducia di un
account recente senza store).


% ─────────────────────────────────────────────────────────────
\section{Analisi NLP del Sentiment}
\label{sec:seller-sentiment}
% ─────────────────────────────────────────────────────────────

Il modulo \texttt{app/services/nlp\_sentiment.py} implementa un sistema
di analisi del sentiment a più livelli, con degradazione graceful
in caso di indisponibilità delle risorse computazionali.

\subsection{Approccio Principale: Classificazione Multilingue con BERT for Sentiment}
\label{subsec:sentiment-bert}

La strategia principale utilizza il modello
\textbf{nlptown/bert-base-multilingual-uncased-sentiment}, un classificatore
basato su BERT caricato tramite la libreria \texttt{transformers} (Hugging Face).
Il modello è specializzato nel predire un rating (da 1 a 5 stelle) a partire da
un testo in ingresso.
Il modello è condiviso nell'applicazione tramite un singleton thread-safe
(\texttt{app/services/model\_singleton.py}), caricato una sola volta
alla prima invocazione e riutilizzato da tutti i componenti.

A differenza del precedente approccio basato su anchor vector e similarità
cosinica, il modello corrente esegue una classificazione diretta in 5 classi.
Il label predetto (\texttt{"1 star"}, \dots, \texttt{"5 stars"}) viene poi
mappato in un punteggio numerico nell'intervallo $[0, 1]$ secondo la seguente
tabella:

\begin{table}[htbp]
\centering
\caption{Mapping predizione BERT → Sentiment Score ($s_{nlp}$)}
\label{tab:bert-score-map}
\begin{tabular}{cc}
\hline
\textbf{BERT Label} & \textbf{Score ($s_{nlp}$)} \\
\hline
5 stars & 0.95 \\
4 stars & 0.75 \\
3 stars & 0.50 \\
2 stars & 0.25 \\
1 star  & 0.05 \\
\hline
\end{tabular}
\end{table}

\subsection{Blended Sentiment: Connessione tra Testo e Rating}
\label{subsec:sentiment-blending}

Per massimizzare la precisione e la robustezza del segnale, il sistema non si
affida esclusivamente all'analisi NLP del testo, ma integra il \textbf{rating
esplicito} (\texttt{Positive}, \texttt{Neutral}, \texttt{Negative}) lasciato
dall'acquirente su eBay per quel singolo feedback.

Il punteggio di sentiment finale per un singolo feedback ($s_i$) è calcolato
come una media pesata tra il valore esplicito normalizzato ($v_{exp}$) e
il punteggio NLP ($s_{nlp}$):

\[
    s_i = 0.6 \cdot v_{exp} + 0.4 \cdot s_{nlp}
\]

dove $v_{exp}$ vale 0.90 per feedback positivi, 0.50 per quelli neutri e 0.10
per quelli negativi. Questa tecnica di \textbf{blending} assicura che il
giudizio testuale agisca come modulatore fine del voto principale, evitando
che ambiguità nel linguaggio (es. sarcasmo non rilevato) portino a stime
eccessivamente errate.

Il sentiment score complessivo del venditore $S$ (usato poi nel trust score)
è la media aritmetica degli $s_i$ calcolati su tutti i feedback analizzati.

\begin{figure}[htbp]
    \centering
    \fbox{\parbox{0.78\linewidth}{\centering\vspace{2cm}
        \textit{[Diagramma: Testo feedback + Rating eBay → nlptown-BERT →
        Predizione Stelle (1-5) → Mapping decimale →
        Blended Score (60/40) → $s_i \in [0,1]$ Media]}\\
        \vspace{2cm}}}
    \caption{Pipeline di calcolo del sentiment via BERT classification e logic di blending.}
    \label{fig:sentiment-pipeline}
\end{figure}

\subsection{Fallback Euristico}
\label{subsec:sentiment-heuristic}

Se il modello di embedding non è disponibile (errore di caricamento,
memoria insufficiente), o se il batch di testi da analizzare è piccolo
($\leq 8$ elementi e \texttt{use\_fast\_path=True}), il sistema degrada
a un approccio \textbf{euristico basato su keyword}:

\begin{enumerate}
    \item Conta le occorrenze di parole positivamente connotate
        (\texttt{\_POS\_WORDS}) e negativamente connotate
        (\texttt{\_NEG\_WORDS}) nel testo.
    \item Calcola il ratio $\frac{\text{pos} - \text{neg}}{\text{hits}+1}$.
    \item Applica la stessa funzione sigmoide del percorso principale.
\end{enumerate}

Questo garantisce che la pipeline produca sempre un output numerico,
indipendentemente dallo stato del sottosistema ML.

\subsection{Labeling LLM per Ingestion RAG}
\label{subsec:sentiment-llm-label}

Per la fase di ingestion nel vettorstore RAG, ogni singolo feedback
viene etichettato con un label discreto (\texttt{POSITIVE},
\texttt{NEGATIVE}, \texttt{NEUTRAL}) tramite una chiamata LLM
(\texttt{extract\_sentiment\_label()}). Anche qui è presente un
fallback euristico se la chiamata LLM fallisce.

\subsection{Cache a Due Livelli}
\label{subsec:sentiment-cache}

Per evitare di ricalcolare l'embedding dello stesso testo più volte
(i feedback di un venditore popolare possono ripresentarsi in richieste
successive), il sistema mantiene una cache in memoria:

\begin{itemize}
    \item Struttura: dizionario \texttt{\{testo: score\}}.
    \item Capacità massima: 5.000 entry.
    \item Eviction: al raggiungimento del limite, il 10\% delle entry
        più vecchie viene rimosso (FIFO approximation).
\end{itemize}

Questa cache è complementare alla cache Redis a livello di pipeline:
Redis evita di ripetere l'intera analisi per un venditore già
valutato di recente; la cache in-memory evita di ricalcolare
l'embedding per singoli testi già visti nella stessa sessione.


% ─────────────────────────────────────────────────────────────
\section{Calcolo del Trust Score}
\label{sec:seller-trust}
% ─────────────────────────────────────────────────────────────

Il modulo \texttt{app/services/trust.py} implementa il calcolo del
\textbf{trust score} composito $T \in [0, 1]$, che sintetizza in un
unico indice la reputazione complessiva del venditore. La formula
finale pesa quattro componenti distinte:

\[
    T = 0.45 \cdot R + 0.25 \cdot S + 0.20 \cdot \text{REC} + 0.10 \cdot C
\]

dove:
\begin{itemize}
    \item $R$ — \textbf{Feedback Ratio}: media dei valori normalizzati
        $\frac{r_i - 1}{4}$ su tutti i feedback con rating valido.
    \item $S$ — \textbf{Sentiment Score}: output della pipeline NLP
        descritta nella sezione precedente.
    \item $\text{REC}$ — \textbf{Recency Score}: media pesata
        temporalmente dei rating, dove il peso di ciascun feedback
        decade esponenzialmente con l'età.
    \item $C$ — \textbf{Confidence}: funzione di saturazione che vale
        0 con zero feedback e converge a 1 con molti feedback.
\end{itemize}

\subsection{Feedback Ratio}
\label{subsec:trust-ratio}

Il feedback ratio $R$ è la media aritmetica dei rating normalizzati.
La normalizzazione lineare
$v = \frac{r - 1}{r_{\max} - 1}$
mappa il minimo eBay (1 = Negative) a 0 e il massimo (5 = Positive)
a 1, con Neutral (3) mappato a 0.5. Rating fuori range o non numerici
vengono esclusi senza generare errori.

\subsection{Recency Score con Decadimento Esponenziale}
\label{subsec:trust-recency}

Un venditore con ottimi feedback di tre anni fa ma pessime recensioni
recenti è più rischioso di quanto il ratio grezzo suggerisca. Il
\textbf{recency score} risolve questo problema assegnando a ogni
feedback un peso inversamente proporzionale alla sua età:

\[
    w_i = e^{-\frac{d_i}{365}}
\]

dove $d_i$ è l'età del feedback in giorni. Con questa funzione, un
feedback di 365 giorni fa vale $e^{-1} \approx 0.37$ volte uno di
oggi; uno di due anni fa vale $e^{-2} \approx 0.14$. Il recency
score è la media pesata:

\[
    \text{REC} = \frac{\sum_i v_i \cdot w_i}{\sum_i w_i}
\]

\subsection{Confidence con Saturazione Esponenziale}
\label{subsec:trust-confidence}

Un venditore con tre feedback positivi è meno affidabile di uno con
trecento. Il fattore di \textbf{confidence} $C$ scala il contributo
del volume dei dati disponibili:

\[
    C = 1 - e^{-n / 50}
\]

Con $n = 50$ feedback, $C \approx 0.63$; con $n = 150$, $C \approx 0.95$.
La costante 50 è calibrata sul dominio eBay: oltre questa soglia, un
venditore ha un campione statisticamente significativo.

\begin{figure}[htbp]
    \centering
    \fbox{\parbox{0.82\linewidth}{\centering\vspace{2cm}
        \textit{[Grafico: curva $C = 1 - e^{-n/50}$ per $n \in [0, 300]$
        — saturazione a 1 dopo circa 150 feedback]}\\
        \vspace{2cm}}}
    \caption{Curva di saturazione del fattore Confidence in funzione
    del numero di feedback disponibili.}
    \label{fig:trust-confidence-curve}
\end{figure}

\subsection{Pesi del Trust Score e Razionale}
\label{subsec:trust-weights}

I pesi della formula composita riflettono una precisa gerarchia
di affidabilità:

\begin{table}[htbp]
\centering
\caption{Componenti del Trust Score e relative motivazioni}
\label{tab:trust-weights}
\begin{tabular}{clcp{0.42\linewidth}}
\hline
\textbf{Peso} & \textbf{Componente} & \textbf{Simbolo} & \textbf{Motivazione} \\
\hline
0.45 & Feedback Ratio   & $R$            & Segnale primario: voto esplicito lasciato dall'acquirente. \\
0.25 & Sentiment Score  & $S$            & Cattura la \emph{qualità} del testo anche a prescindere dalla stella. \\
0.20 & Recency Score    & $\text{REC}$   & Penalizza i venditori con performance recenti deteriorate. \\
0.10 & Confidence       & $C$            & Corregge l'ottimismo su campioni piccoli. \\
\hline
\end{tabular}
\end{table}

Il sentiment score (25\%) è intenzionalmente più basso del feedback
ratio (45\%) perché i feedback eBay sono già strutturalmente orientati:
la maggior parte degli acquirenti lascia solo voti positivi, rendendo
il testo il segnale più informativo per distinguere i casi intermedi
e negativi.


% ─────────────────────────────────────────────────────────────
\section{AI Summary Generato da LLM}
\label{sec:seller-ai-summary}
% ─────────────────────────────────────────────────────────────

Come passo finale della pipeline, viene generato un \textbf{sommario
testuale} in italiano del profilo del venditore, direttamente da LLM.
Fino a 100 testi di feedback vengono campionati, concatenati e passati
a \texttt{call\_llm()} con un prompt istruttivo:

\begin{quotation}
\textit{``Analizza queste recenti e autentiche recensioni di acquirenti
sul venditore '\{seller\_name\}'. Scrivi in lingua italiana al massimo
2 o 3 frasi (max 50 parole) che riassumono il giudizio generale, i
pregi e i potenziali difetti o rischi emersi. Sii conciso e diretto,
senza convenevoli.''}
\end{quotation}

Il risultato è un campo \texttt{ai\_summary} nel payload di risposta,
pensato per essere mostrato direttamente nella UI senza ulteriore
elaborazione. Se la chiamata LLM fallisce, il campo è semplicemente
assente (\texttt{None}), e la pipeline non solleva eccezioni.

Questa strategia combina il vantaggio della sintesi automatica
(accessibile all'utente non tecnico) con la robustezza numerica
del trust score (usabile dall'agente per ranking e confronti).


% ─────────────────────────────────────────────────────────────
\section{Caching Redis e Politica di Invalidazione}
\label{sec:seller-cache}
% ─────────────────────────────────────────────────────────────

La pipeline implementa una cache Redis con TTL di 5 minuti. La chiave
di cache codifica il venditore e il volume di feedback richiesto:

\begin{verbatim}
cache_key = f"seller_analysis:{seller_name.lower()}:{needed}"
\end{verbatim}

dove \texttt{needed = page * limit}. Questa strategia garantisce che
richieste con parametri di paginazione diversi abbiano chiavi distinte,
evitando la restituzione di risultati troncati per richieste con
\texttt{limit} maggiore.

Il TTL di 5 minuti è un compromesso deliberato: abbastanza breve da
non presentare dati obsoleti in scenari di disputa o feedback
recentemente aggiunti, abbastanza lungo da assorbire picchi di richieste
parallele per lo stesso venditore (tipico quando più utenti cercano
lo stesso prodotto nello stesso momento).

In assenza di Redis, la classe \texttt{redis\_client} degrada
silenziosamente a un dizionario in-memory, preservando la funzionalità.


% ─────────────────────────────────────────────────────────────
\section{Normalizzazione dell'Output}
\label{sec:seller-normalizer}
% ─────────────────────────────────────────────────────────────

Il risultato grezzo della pipeline viene elaborato dal normalizer
\texttt{\_normalize\_seller\_output()} (\texttt{app/mcp/normalizers.py}),
che aggiunge un campo \texttt{summary} leggibile in linguaggio naturale,
basato su soglie semantiche condivise:

\begin{table}[htbp]
\centering
\caption{Soglie e label semantiche del normalizer}
\label{tab:normalizer-thresholds}
\begin{tabular}{lll}
\hline
\textbf{Metrica} & \textbf{Soglia} & \textbf{Label} \\
\hline
Trust Score & $\geq 0.85$ & ``eccellente'' \\
Trust Score & $\geq 0.70$ & ``affidabile'' \\
Trust Score & $< 0.70$    & ``scarsa affidabilità'' \\
\hline
Sentiment Score & $\geq 0.70$ & ``positivi'' \\
Sentiment Score & $\leq 0.40$ & ``negativi'' \\
Sentiment Score & intermedio  & ``misti'' \\
\hline
\end{tabular}
\end{table}

Un esempio di summary generato dal normalizer:

\begin{quotation}
\textit{``Analisi completata per tech-store-it (287 feedback). Profilo
eccellente (Trust: 88\%), feedback positivi (Sentiment: 76\%).''}
\end{quotation}

Questo campo è consumato direttamente dall'agente ReAct per formulare
la risposta in linguaggio naturale all'utente.


% ─────────────────────────────────────────────────────────────
\section{Integrazione con la Search Pipeline e il Reranker LTR}
\label{sec:seller-ltr-integration}
% ─────────────────────────────────────────────────────────────

Il trust score prodotto dalla Seller Pipeline non è confinato al tool
\texttt{analyze\_seller}: viene usato come \textbf{feature di ranking}
nella pipeline di ricerca prodotti. Il modello LTR
(\emph{Learning to Rank}) incorpora il trust score del venditore tra
le sue feature, consentendo al reranker di penalizzare automaticamente
i prodotti venduti da account con bassa reputazione, anche quando il
prezzo e la pertinenza semantica sono competitivi.

Questa integrazione trasforma il trust score da semplice etichetta
descrittiva a componente attivo della strategia di ottimizzazione
dei risultati, chiudendo il ciclo tra analisi della reputazione e
qualità delle raccomandazioni.


% ─────────────────────────────────────────────────────────────
\section{Riepilogo delle Strategie Implementate}
\label{sec:seller-summary}
% ─────────────────────────────────────────────────────────────

La Seller Pipeline combina un insieme articolato di strategie tecniche,
ciascuna indirizzata a un aspetto specifico del problema:

\begin{table}[htbp]
\centering
\caption{Riepilogo delle strategie implementate nella Seller Pipeline}
\label{tab:seller-strategies}
\begin{tabular}{p{0.28\linewidth}p{0.62\linewidth}}
\hline
\textbf{Strategia} & \textbf{Descrizione} \\
\hline
Fetch multi-pagina parallelo & Tutte le pagine eBay necessarie vengono richieste in simultanea con \texttt{asyncio.gather()}, riducendo la latenza di acquisizione. \\
Fetch dati eterogenei in parallelo & Feedback, dati utente e dati store vengono acquisiti con tre chiamate API concorrenti. \\
Classificazione sentiment con BERT & Il modello \texttt{nlptown/bert-base-multilingual-uncased-sentiment} opera su testi multilingue per classificare il sentiment in stelle (1-5). \\
Blended scoring & Il punteggio NLP viene combinato con il rating esplicito di eBay (60\% rating, 40\% NLP) per una valutazione più solida. \\
Trust score composito & Formula ponderata con quattro segnali (ratio, sentiment, recency, confidence) calibrati sul dominio e-commerce. \\
Decadimento esponenziale temporale & I feedback recenti pesano proporzionalmente di più nella valutazione della reputazione. \\
Confidence adattiva & Il contributo dei dati scala con il volume disponibile, evitando stime ottimistiche su campioni piccoli. \\
AI Summary tramite LLM & Sintesi in italiano del profilo venditore generata dall'LLM a partire dai feedback reali. \\
Cache a due livelli & Redis (5 min, livello pipeline) + dizionario in-memory (5.000 entry, livello testo) per evitare ricalcoli ridondanti. \\
Normalizzazione semantica dell'output & Label human-readable generati da soglie condivise, direttamente fruibili dall'agente. \\
Integrazione con LTR Reranker & Il trust score è feature del reranker di ricerca, impattando il ranking dei prodotti. \\
\hline
\end{tabular}
\end{table}

Il capitolo successivo descrive la pipeline di ricerca prodotti
(\S\ref{chap:search}), nella quale il trust score del venditore
si integra con le feature semantiche e di prezzo per il ranking
finale degli articoli.
